\talkslide{attacker-loss}{Learning the attacker: online DPO}{Timing: provisional}{For the attacker, the candidates are two prompts derived from the same seed $x_s$ (written $s$ in the paper). The preferred and rejected attacks are $x_w$ and $x_l$; the preceding example has $x_w=x_2$ and $x_l=x_1$.}{
\dpolossview{0}
}
\speak{The diagram shows a schematic reference-relative score profile, with the preferred and rejected candidates as individual points. Language-model outputs are discrete and the candidate ordering is arbitrary. The connecting curve is only a visual guide, not a probability density or measured curve. Scores can be positive or negative; the schematic has no numeric vertical scale. The vertical arrows depict raising the preferred candidate’s score and lowering the rejected candidate’s score. The exact objective compares their reference-relative log probabilities, not their absolute probabilities alone. The score is beta times the log probability relative to the reference model, not a scalar score emitted by the judge. The expectation averages this loss over freshly collected preference pairs. Winner and loser are denoted w and l. The compact notation suppresses the fixed conditioning variables inside the score. For the attacker, the candidates are two prompts derived from the same seed $x_s$ (written $s$ in the paper). The preferred and rejected attacks are $x_w$ and $x_l$; the preceding example has $x_w=x_2$ and $x_l=x_1$. Online DPO turns a preference between these attacks into a likelihood-ratio update relative to the reference attacker. Here w and l denote winner and loser, phi denotes attacker parameters, beta scales the update, and sigma is the sigmoid. New attacks and judgments are collected throughout training. In the full AdvGame rule, faithfulness comes first, and each attack is assessed through its preferred defender response. That optimistic comparison makes the attacker challenge a better response rather than exploiting a worst sample. We will show the exact faithfulness gate and rollout tree in the defense section. The defender later uses the corresponding loss on response pairs. Source detail: AdvGame, Eq. (9) and Algorithm 1.}
